\documentclass[10pt,a4paper]{amsart}
\usepackage[margin=19mm]{geometry}
\usepackage{amsmath,amssymb,mathtools}
\usepackage[scr=rsfs,cal=euler]{mathalfa}
\usepackage{xfrac}
\usepackage[unicode,hidelinks,bookmarks=false]{hyperref}
\newcommand{\ie}{i.e.\ }
\newcommand{\cf}{cf.\ }
\newcommand{\resp}{respectively\ }
\newcommand{\Subsection}{\subsection}
\providecommand{\nobreakdash}{\nobreak\hbox{-}}
\setlength{\emergencystretch}{2em}
\multlinegap0pt
\usepackage{bm}
\usepackage{bbm}
\usepackage{dsfont}
\usepackage{xspace}
\usepackage{cancel}
\usepackage{mathtools}
\usepackage{amsthm}
\makeatletter
\@ifundefined{cor}{\newtheorem{cor}{Corollary}}{}
\@ifundefined{keylem}{\newtheorem{keylem}{Key-Lemma}}{}
\makeatother
\def\sH{\mathscr{H}}
\def\sP{\mathscr{P}}
\def\sQ{\mathscr{Q}}
\def\sS{\mathscr{S}}
\def\sV{\mathscr{V}}
\def\sZ{\mathscr{Z}}
\title[Optimal Control of Pseudo-Parabolic KWC Systems for Grain Bo]{Optimal Control of Pseudo-Parabolic KWC Systems for Grain Boundary Motion}
\author{Harbir Antil, Daiki Mizuno, Ken Shirakawa}
\date{Source: arXiv:2506.09407v1 (CC BY 4.0)}
\begin{document}
\maketitle

\noindent\emph{Source and licence.} Optimal Control of Pseudo-Parabolic KWC Systems for Grain Boundary Motion. Version arXiv:2506.09407v1 (CC BY 4.0), distributed under the Creative Commons Attribution 4.0 International licence, as recorded in the arXiv licence metadata for this exact version and shown on the abstract page https://arxiv.org/abs/2506.09407v1. Full rights notice: licenses/arxiv-2506.09407-CC-BY-4.0.txt.\par\smallskip

\noindent\textit{Editorial scope.} Counted unit: the Cor whose source label is \texttt{cor\_bddness} in the article (source0.tex lines 1662--1676, source label \texttt{cor\_bddness}), delivered together with the article's own complete original proof of it (source0.tex lines 1677--1749, one \texttt{proof} environment of 73 lines). No mathematics is written, repaired or re-derived: statement and proof are the article's own text line for line. Required context retained uncounted: lem:P\_est (lines 1095--1122); cor:P\_est (lines 1515--1530). Every definition, display and label of those lines is retained unchanged. Omitted from this package and not redistributed: 1--1094, 1123--1514, 1531--1661, 1750--2512. Nothing inside a retained excerpt is omitted or altered: every retained body is delivered exactly as the article's lines are written, so no omission receipt of any kind accompanies this package. Citations: the retained excerpts carry no citation command at all, so no bibliography entry of the article is transcribed and no reference is redistributed. Reference adaptation: none was needed and none was made; every reference inside the delivered unit resolves inside it, so no unresolved reference or citation is produced. Printed slips kept literal: no printed word, symbol, hypothesis or number of the counted statement or of its proof was altered. Layout and class: the article's own class is \texttt{sinst} with packages including amsbsy, amsfonts, amsmath, amssymb, amsthm, array, bm, color, enumerate, mathrsfs, latexsym, hyperref, mathtools, todonotes; the wrapper is the lane's pinned \texttt{amsart} editorial wrapper at 10pt on A4, which restates the theorem-like environments the retained text needs and adds the article's own macro definitions that the retained text needs (\texttt{\textbackslash sH}, \texttt{\textbackslash sP}, \texttt{\textbackslash sQ}, \texttt{\textbackslash sS}, \texttt{\textbackslash sV}, \texttt{\textbackslash sZ}), with extra wrapper packages (bm, bbm, dsfont, xspace, cancel, mathtools). The delivered numbering is the unit's own. Assets: no retained span contains a figure environment, a graphics-inclusion command, a PDF-page inclusion, a table or a TikZ picture; no image file is copied into this package, the package declares no asset dependency and no image file is redistributed with it. Licence: version arXiv:2506.09407v1 of this article is distributed under the Creative Commons Attribution 4.0 International licence, as recorded in the arXiv licence metadata for this exact version and shown on the abstract page https://arxiv.org/abs/2506.09407v1. A full rights notice is delivered at licenses/arxiv-2506.09407-CC-BY-4.0.txt. Characters: every retained excerpt is pure ASCII, so the delivered engine, XeLaTeX with its default Latin Modern text font, reports no missing character.
\begin{keylem} \label{lem:P_est}
  For any $\tau \in (0,\tau_1)$, our time-discretization scheme (TD)$_\tau$ admits a unique solution. Moreover, There exists a small time-step size $\tau_2 \in (0,\tau_1)$ such that the following estimates (i) and (ii) hold.
  
  \smallskip
  \noindent
  (i) Let $C_1 = C_1(a,b,c,\lambda,\xi,\omega)$ be a positive constant which depends on $|\partial_t a|_\sH$, $|b|_\sH$, $|c|_\sH$, $|\lambda|_\sH$, $|\xi|_\sH$, and $|\omega|_{[L^\infty(Q)]^N}$, given as:
  \begin{gather}
    C_1 := \frac{8((C_V^{L^4})^2 + 1)}{1 \wedge \delta_a \wedge \mu^2 \wedge \nu^2}\bigl( |\partial_t a|_\sH + |b|_\sH + |c|_\sH + |\lambda|_\sH + |\xi|_\sH + |\omega|_{[L^\infty(Q)]^N} + 1\bigr).
  \end{gather}
  Then, it holds that:
  \begin{align}
    &|p_i|_H^2 + \mu^2 |\nabla p_i|_{[H]^N}^2 + |\sqrt{a_i}z_i|_H^2 + \nu^2 |\nabla z_i|_{[H]^N}^2 
    \\
    &\quad \leq e^{4 C_1 (T + 1)} \bigl( |p_0|_H^2 + \mu^2 |\nabla p_0|_{[H]^N}^2 + |\sqrt{a(0)}z_0|_H^2 + \nu^2 |\nabla z_0|_{[H]^N}^2 \bigr)
    \\
    &\quad\qquad + 2(|h|_{\sV^*}^2 + |k|_{\sV^*}^2), \mbox{ for any } i = 1,2,3,\dots, n_\tau. \label{P_TD_est1}
  \end{align}
  (ii) There exists a positive constant $C_2 = C_2(a) > 0$, depending on $\delta_a$ and $|a(0)|_H$ such that
  \begin{gather}
    \frac{1 \wedge \mu^2}{4\tau}\sum_{i=1}^{n_\tau} |p_i - p_{i-1}|_V^2 \leq C_2 e^{5C_1 (T + 1)} (|p_0|_V^2 + |z_0|_V^2 + |h|_{\sV^*}^2 + |k|_{\sV^*}^2), \label{P_TD_est2}
  \end{gather}
  and
  \begin{align}
    &\frac{\delta_a \wedge \nu^2}{4\tau}\sum_{i=1}^{n_\tau} |z_i - z_{i-1}|_V^2 
    \\
    &\qquad \leq C_2 e^{5 C_1 (T + 1)}(|A|_{[L^\infty(Q)]^{N \times N}}^2 + 1) (|p_0|_V^2 + |z_0|_V^2 + |h|_{\sV^*}^2 + |k|_{\sV^*}^2). \label{P_TD_est3}
  \end{align}
\end{keylem}

\begin{cor}\label{cor:P_est}
  Let $[p,z]$ be the solution to the system (P) corresponding to a septuplet $[a,b,c,\lambda,\xi,\omega,A] \in \sS$, initial data $[p_0,z_0] \in [V]^2$ and forcing pair $[h,k] \in [\sV^*]^2$. Then, the solution $[p,z]$ fulfills the following estimates:
    \begin{align}
      |p|_{C([0,T];V)}^2 + |z|_{C([0,T];V)}^2 &\leq \frac{1 + \mu^2 + \nu^2 + (C_V^{L^4})^2|a(0)|_H}{1 \wedge \delta_a \wedge \mu^2 \wedge \nu^2}e^{4 C_1 (T+1)}\cdot
      \\
      &\qquad \cdot (|p_0|_V^2 + |z_0|_V^2 + |h|_{\sV^*}^2 + |k|_{\sV^*}^2), \label{P_est_97}
    \end{align}
    \begin{gather}
      \frac{1 \wedge \mu^2}{4}|\partial_t p|_{\sV}^2 \leq C_2 e^{5C_1 (T + 1)} (|p_0|_V^2 + |z_0|_V^2 + |h|_{\sV^*}^2 + |k|_{\sV^*}^2), \label{P_TD_est98}
    \end{gather}
    and
    \begin{align}
      &\frac{\delta_a \wedge \nu^2}{4}|\partial_t z|_{\sV}^2 \leq C_2 e^{5 C_1 (T + 1)}(|A|_{[L^\infty(Q)]^{N \times N}}^2 + 1) (|p_0|_V^2 + |z_0|_V^2 + |h|_{\sV^*}^2 + |k|_{\sV^*}^2), \label{P_TD_est99}
    \end{align}
    where the constants $C_1, C_2$ are as in Key-Lemma \ref{lem:P_est}.
\end{cor} 

\begin{cor}[Linearity and boundedness of solution operator] \label{cor_bddness}
  Let us denote 
  \\
  $\sZ = [W^{1,2}(0,T;V)]^2$, and let $\sZ$ be a Banach space endowed with a norm $\| \cdot \|: \sZ \longrightarrow [0,\infty)$, defined as:
    \begin{equation}
      \| [p,z] \| := |[p,z]|_{[W^{1,2}(0,T;V)]^2} + |[p,z]|_{C([0,T];[V]^2)}.
    \end{equation}
    We let $[a,b,c,\lambda,\xi,\omega,A] \in \sS$, $[p_0, z_0] \in [V]^2$, $[h,k] \in [\sV^*]^2$, and $[p,z] := \sP([p_0,z_0],[h,k])$ $= \sP(a,b,c,\lambda,\xi,\omega,A)([p_0,z_0],[h,k])$. Then, there exist two positive constants $M_0^* = M_0^*(a,b,c,\lambda,\xi,\omega,A)$, $M_1^* = M_1^*(a,b,c,\lambda,\xi,\omega,A)$ such that:
    \begin{gather}
      M_0^* |[[p_0,z_0],[h,k]]|_{[V]^2 \times [\sV^*]^2} \leq \|[p,z]\| \leq M_1^* |[[p_0,z_0],[h,k]]|_{[V]^2 \times [\sV^*]^2}, 
      \\
      \mbox{ for any } [p_0, z_0] \in [V]^2, \mbox{ and } [h,k] \in [\sV^*]^2. \label{sP_bdd}
    \end{gather}
    Moreover, $\sP$ is an bijective operator from $[V]^2 \times [\sV^*]^2$ to $\sZ$, Therefore, $\sP$ is an isomorphism between a Hilbert space $[V]^2 \times [\sV^*]^2$ and a Banach space $(\sZ, \| \cdot \|)$.
\end{cor}

\begin{proof}
  Corollary \ref{cor:P_est} implies that $\sP$ is a linear bounded operator from $[V]^2 \times [\sV^*]^2$ to $(\sZ, \| \cdot \|)$ with the constant:
  \begin{equation}
    M_1^* := \left(\frac{48(1 + \mu^2 + \nu^2 + (C_V^{L^4})^2)|a(0)|_H}{1 \wedge \delta_a \wedge \mu^2 \wedge \nu^2} C_2 e^{5 C_1 (T+1)}(|A|_{[L^\infty(Q)]^{N \times N}} + 1)\right)^\frac{1}{2}.
  \end{equation}

  Now, we define a linear operator $\sQ = \sQ(a,b,c,\lambda,\xi,\omega,A): \sZ \longrightarrow [\sV^*]^2$ which maps $[p,z]$ to a pair of bounded linear functionals $[\tilde h, \tilde k] =: \sQ(p,z) \in [\sV^*]^2$, given as:
  \begin{gather}
    \langle \tilde h, \varphi \rangle_\sV := \int_0^T (\partial_t p(t) + \omega(t) \cdot \nabla z(t), \varphi(t))_H \,dt + \int_0^T \int_\Omega (\lambda(t) p(t) + \xi(t)z(t)) \varphi(t) \,dxdt 
    \\
    + \int_0^T (\nabla (p + \mu^2 \partial_t p)(t), \nabla \varphi(t))_{[H]^N} \,dt, \mbox{ for any } \varphi \in \sV,
  \end{gather}
  and
  \begin{gather}
    \langle \tilde k,\psi \rangle_\sV := \int_0^T \int_\Omega (a(t) \partial_t z(t) + b(t) z(t) + c(t) p(t)) \psi(t) \,dxdt 
    \\
    + \int_0^T (A(t) \nabla z(t) + \nu^2\nabla \partial_t z(t) + p(t)\omega(t), \nabla \psi(t))_{[H]^N} \,dt, \mbox{ for any } \psi \in \sV.
  \end{gather}
  Then, one can see that $\sQ$ is a linear bounded operator from $\sZ$ to $[\sV^*]^2$ together with the following computations:
  \begin{align}
    &\quad |\langle \tilde h, \varphi \rangle_\sV| 
    \\
    &\leq (|\partial_t p|_\sH + |\omega|_{[L^\infty(Q)]^N} |\nabla z|_{[\sH]^N})|\varphi|_\sH + (C_V^{L^4})^2(|\lambda|_\sH |p|_{C([0,T];V)} + |\xi|_\sH |z|_{C([0,T];V)}) |\varphi|_\sV
    \\
    &\qquad + (|\nabla p|_{[\sH]^N} + \mu^2 |\nabla \partial_t p|_{[\sH]^N}) |\nabla \varphi|_{[\sH]^N} 
    \\
    &\leq 2 (1 + \mu^2)\bigl( (C_V^{L^4})^2 (|\lambda|_\sH + |\xi|_\sH) + |\omega|_{[L^\infty(Q)]^N} + 1\bigr)\| [p,z] \| \cdot |\varphi|_\sV,
    \\
    &\qquad\qquad\qquad\qquad\qquad\qquad \mbox{ for any } \varphi \in \sV,
  \end{align}
  and
  \begin{align}
    &\quad|\langle \tilde k,\psi \rangle_\sV| 
    \\
    &\leq (C_V^{L^4})^2(|a|_{C([0,T];H)} |\partial_t z|_\sV + |b|_\sH |z|_{C([0,T];V)} + |c|_\sH |p|_{C([0,T];V)})|\psi|_\sV 
    \\
    &\qquad + (|A|_{[L^\infty(Q)]^{N \times N}} |\nabla z|_{[\sH]^N} + \nu^2 |\nabla \partial_t z|_{[\sH]^N} + |\omega|_{[L^\infty(Q)]^N} |p|_\sH) |\nabla \psi|_{[\sH]^N}
    \\
    &\leq 2 (1 + \nu^2) \bigl( (C_V^{L^4})^2 (|a|_{C([0,T];H)} + |b|_\sH + |c|_\sH) + |A|_{[L^\infty(Q)]^{N \times N}} + |\omega|_{[L^\infty(Q)]^N} + 1\bigr)\cdot 
    \\
    &\qquad \cdot \| [p,z] \| \cdot |\psi|_\sV.
  \end{align}
  and hence,
  \begin{align}
    &\quad |\sQ(p,z)|_{[\sV^*]^2} \leq \sqrt{2}(|\tilde h|_{\sV^*} + |\tilde k|_{\sV^*}) \label{P_bdd_1}
    \\
    &\leq \tilde C_4 (|a|_{C([0,T];H)} + |b|_\sH + |c|_\sH + |\lambda|_\sH + |\xi|_\sH + |A|_{[L^\infty(Q)]^{N \times N}} + |\omega|_{[L^\infty(Q)]^N} + 1) \cdot
    \\
    &\qquad \cdot \| [p,z] \|,
  \end{align}
  with
  \begin{equation}
    \tilde C_4 := 4\sqrt{2} (1 + \mu^2 + \nu^2)(1 + (C_V^{L^4})^2).
  \end{equation}

  Besides, we define a linear operator $\overline{\sP} = \overline{\sP}(a,b,c,\lambda,\xi,\omega,A): \sZ \longrightarrow [V]^2 \times [\sV^*]^2$, by letting:
  \begin{equation}
    \overline{\sP}: [p,z] \in \sZ \mapsto \overline{\sP}([p,z]) := \Bigl( [p(0), z(0)], \sQ(p,z) \Bigr) \in [V]^2 \times [\sV^*]^2.
  \end{equation}
  Then, we can check that $\overline{\sP}$ is the inverse map of $\sP$. Moreover, noting that:
  \begin{equation}
    |[p(0), z(0)]|_{[V]^2} \leq |[p,z]|_{C([0,T];[V]^2)} \leq \|[p,z]\|. \label{P_bdd_2}
  \end{equation}
  we can say that $\overline{\sP}$ is a bounded linear operator from $\sZ$ to $[V]^2 \times [\sV^*]^2$, asserted in \eqref{sP_bdd}, with the constant $M_0^*$, given by:
  \begin{align}
    &M_0^* = 2^{-\frac{1}{2}}(\tilde C_4 + 1)^{-1} \cdot
    \\
    &\cdot (|a|_{C([0,T];H)} + |b|_\sH + |c|_\sH + |\lambda|_\sH + |\xi|_\sH + |A|_{[L^\infty(Q)]^{N \times N}} + |\omega|_{[L^\infty(Q)]^N} + 1)^{-1}.
  \end{align}

  Thus, we finish the proof of Corollary \ref{cor_bddness}.

\end{proof}

\end{document}
